\documentclass[11pt]{scrartcl}

\usepackage[stixtwo]{fontsetup}
\usepackage{amsmath}
\usepackage{tikz}

\begin{document}
\begin{enumerate}
	\item \[
		      z
		      = \frac{1 - 5i}{1 - i}
		      = \frac{(1 - 5i)(1 + i)}{(1 - i)(1 + i)}
		      = \frac{1 - 5i + i - 5i^2}{1 - i^2}
		      = \frac{6 - 4i}{2}
		      = 3 - 2i,
	      \]
	      tehát $\operatorname{Re} z = 3$ és $\operatorname{Im} z = -2$.

	      $|1 + \sqrt{3} i| = \sqrt{1 + 3} = 2$, és mivel
	      $\operatorname{Im}(1 + \sqrt{3} i) = \sqrt{3} > 0$, ezért
	      \(
	      \arg(1 + \sqrt{3} i)
	      = \arccos \frac{\operatorname{Re}(1 + \sqrt{3} i)}{|1 + \sqrt{3} i|}
	      = \arccos \frac{1}{2}
	      = 60^\circ
	      \).
	      A Moivre-képlettel
	      \[
		      w
		      = 2^{2022} \left( \cos(2022 \cdot 60^\circ) + i \sin(2022 \cdot 60^\circ) \right)
		      = 2^{2022} \left( \cos 0^\circ + i \sin 0^\circ \right)
		      = 2^{2022},
	      \]
	      mert $2022 \cdot 60^\circ = 337 \cdot 360^\circ$ (a $2022$ osztható
	      $6$-tal). Tehát $\operatorname{Re} w = 2^{2022}$ és
	      $\operatorname{Im} w = 0$.

	      A $z$ és az $1 + \sqrt{3} i$ szám a komplex számsíkon ($w$ a hossza
	      miatt nem fér az ábrára; a pozitív valós tengelyen van):
	      \begin{center}
		      \begin{tikzpicture}[scale=0.8, >=stealth]
			      \draw[->] (-1, 0) -- (4, 0) node[right] {$\operatorname{Re}$};
			      \draw[->] (0, -2.8) -- (0, 2.4) node[above] {$\operatorname{Im}$};
			      \draw (0, 0) -- (60:2);
			      \draw (0.6, 0) arc[start angle=0, end angle=60, radius=0.6];
			      \node at (30:0.95) {$60^\circ$};
			      \fill (60:2) circle (2pt) node[above right] {$1 + \sqrt{3} i$};
			      \fill (3, -2) circle (2pt) node[below right] {$z = 3 - 2i$};
		      \end{tikzpicture}
	      \end{center}

	\item $z^3 + i z^2 + 2z = z \left( z^2 + i z + 2 \right)$ gyökei $z_1 = 0$ és
	      \[
		      z_{2, 3}
		      = \frac{-i \pm \sqrt{i^2 - 4 \cdot 1 \cdot 2}}{2}
		      = \frac{-i \pm \sqrt{-9}}{2}
		      = \frac{-i \pm 3i}{2},
	      \]
	      tehát $z_2 = i$ és $z_3 = -2i$. A gyökök a komplex számsíkon:
	      \begin{center}
		      \begin{tikzpicture}[scale=0.7, >=stealth]
			      \draw[->] (-1.6, 0) -- (1.6, 0) node[right] {$\operatorname{Re}$};
			      \draw[->] (0, -2.6) -- (0, 1.8) node[above] {$\operatorname{Im}$};
			      \fill (0, 0) circle (2pt) node[below left] {$z_1$};
			      \fill (0, 1) circle (2pt) node[right] {$z_2 = i$};
			      \fill (0, -2) circle (2pt) node[right] {$z_3 = -2i$};
		      \end{tikzpicture}
	      \end{center}
\end{enumerate}
\end{document}
